\documentclass[11pt]{article}
\usepackage[utf8]{inputenc}
\usepackage{amsmath,amssymb}
\usepackage{hyperref}
\title{Efficient Sleep Memory Consolidation: A Resource-Aware Approach}
\author{Anonymous}
\date{}
\begin{document}
\maketitle

\begin{abstract}
This paper studies Sleep Memory Consolidation. Grounded in a real, citable literature \cite{ref1,ref2,ref3,ref4,ref5,ref6,ref7,ref8},
we propose a focused method and report \textbf{illustrative} experiments.
All references below are real and verifiable (OpenAlex/arXiv); the reported
numbers are simulated to illustrate intended behaviour.
\end{abstract}

\section{Introduction}
Sleep Memory Consolidation is an active research area. This work targets a concrete gap identified from
the recent literature and contributes a method together with an evaluation protocol.

\section{Related Work (real literature)}
The following works, retrieved from public bibliographic APIs, frame our study:
\begin{itemize}
\item Diekelmann et al. (2010) — "The memory function of sleep", Nature reviews. Neuroscience (cited 3770).
\item Wilson et al. (1994) — "Reactivation of Hippocampal Ensemble Memories During Sleep", Science (cited 3214).
\item Rasch et al. (2013) — "About Sleep's Role in Memory", Physiological Reviews (cited 2845).
\item Tononi et al. (2014) — "Sleep and the Price of Plasticity: From Synaptic and Cellular Homeostasis to Memory Consolidation and Integration", Neuron (cited 2429).
\item Buzsáki (2015) — "Hippocampal sharp wave‐ripple: A cognitive biomarker for episodic memory and planning", Hippocampus (cited 1899).
\item Stickgold (2005) — "Sleep-dependent memory consolidation", Nature (cited 1865).
\end{itemize}

\section{Method}
We reduce the dominant computational cost via adaptive resource allocation, activating only the components needed per input. We instantiate this idea for Sleep Memory Consolidation and analyse its cost/quality trade-off.

\section{Experiments (Illustrative)}
\begin{center}
\begin{tabular}{lcc}
System & Primary Metric & Efficiency \\ \hline
Strong baseline & 0.812 & 1.00$\times$ \\
Ablation & 0.828 & 0.92$\times$ \\
\textbf{Ours} & \textbf{0.851} & \textbf{0.71$\times$}
\end{tabular}
\end{center}
\emph{Metrics simulated to illustrate the intended effect; not measured.}

\section{Conclusion}
We connected a real literature to a concrete method for Sleep Memory Consolidation. Future work will
validate the illustrative results with real experiments.

\begin{thebibliography}{9}
\bibitem{ref1} Susanne Diekelmann, Jan Born. The memory function of sleep. \emph{Nature reviews. Neuroscience}, 2010. DOI:10.1038/nrn2762
\bibitem{ref2} Matthew A. Wilson, Bruce L. McNaughton. Reactivation of Hippocampal Ensemble Memories During Sleep. \emph{Science}, 1994. DOI:10.1126/science.8036517
\bibitem{ref3} Björn Rasch, Jan Born. About Sleep's Role in Memory. \emph{Physiological Reviews}, 2013. DOI:10.1152/physrev.00032.2012
\bibitem{ref4} Giulio Tononi, Chiara Cirelli. Sleep and the Price of Plasticity: From Synaptic and Cellular Homeostasis to Memory Consolidation and Integration. \emph{Neuron}, 2014. DOI:10.1016/j.neuron.2013.12.025
\bibitem{ref5} György Buzsáki. Hippocampal sharp wave‐ripple: A cognitive biomarker for episodic memory and planning. \emph{Hippocampus}, 2015. DOI:10.1002/hipo.22488
\bibitem{ref6} Robert Stickgold. Sleep-dependent memory consolidation. \emph{Nature}, 2005. DOI:10.1038/nature04286
\bibitem{ref7} Lisa Marshall, Halla Helgadóttir, Matthias Mölle et al.. Boosting slow oscillations during sleep potentiates memory. \emph{Nature}, 2006. DOI:10.1038/nature05278
\bibitem{ref8} Daoyun Ji, Matthew Wilson. Coordinated memory replay in the visual cortex and hippocampus during sleep. \emph{Nature Neuroscience}, 2006. DOI:10.1038/nn1825
\end{thebibliography}
\end{document}
